\subsection{Model Checking}
\begin{definition}[]{Model Checking}
    Model checking is an automated technique that, given
    a finite-state model of a system and a formal property,
    systematically checks whether this property holds for
    (a given state in) that model.
\end{definition}

Model checkers enumerate all possible states of a system, either through explicitly representing state through concrete values or symbolically through (boolean) formulas.

They are primarily used to analyze system \bi{designs}, and not implementations and are often used to analyze deadlocks, the reachability of undesired states and protocol violations.


\subsubsection{The Model Checking Process}
The first and most important phase is the \bi{modeling phase}, where we model the system in the description language of the model checker (here Promela).
It also includes formalizing the properties to be checked in said language.

Next, we run the model checker to check the validity of the model. 
In the case of this course, we use \texttt{spin}, and we can run a promela model using \texttt{spin -x <promela file>.pml},
which wraps \texttt{spin -a <promela file>.pml}, \texttt{gcc <promela file>.c} and \texttt{./a.out} into a single command.

After running, it is time to analyze the output of the model checker. If the property is violated, analyze the found conter example.
If the mdeol is too large, it can happen that the checker runs out of memory. In that case, reduce the model and try again.
